% ============================================================================
% Appendix for PRL_attemption.tex
% Full mathematical details omitted from the 5-page PRL main text.
% ============================================================================
\documentclass[11pt]{article}

\usepackage{amsmath,amssymb,bm,geometry}
\geometry{margin=1in}

\title{Appendix: Detailed Derivations for the 2D Supersonic Marshak-Wave Model}
\author{Yuval Kochavi et al.}
\date{\today}

\begin{document}
\maketitle

\section*{Scope}
This appendix provides the detailed steps that are only summarized in the PRL manuscript: (i) the numerical closure of the 1D Hammer--Rosen model with a Marshak boundary, (ii) the radial eigenvalue problem in cylindrical geometry and its link to albedo, and (iii) the explicit wall-loss and wall-ablation accounting used in the 2D correction.

\section{1D Core Equations and Marshak Closure}
We use the standard power-law material model
\begin{equation}
\kappa^{-1}=gT^{\alpha}\rho^{-\lambda},
\qquad
e=fT^{\beta}\rho^{-\mu},
\end{equation}
and define
\begin{equation}
H(t)\equiv T_s^{4+\alpha}(t),
\qquad
\varepsilon\equiv\frac{\beta}{4+\alpha},
\qquad
C\equiv\frac{16g\sigma_{\mathrm{SB}}}{3f(4+\alpha)}\rho^{2-\mu-\lambda}.
\end{equation}
The Hammer--Rosen front relation is
\begin{equation}
x_F^2(t)=\frac{2+\varepsilon}{1-\varepsilon}C\,H^{-\varepsilon}(t)\int_0^t H(t')\,dt'.
\label{eq:xF_HR_appendix}
\end{equation}
The material energy per unit cross-sectional area is
\begin{equation}
E(t)=f\rho^{1-\mu}(1-\varepsilon)x_F(t)H^{\varepsilon}(t),
\label{eq:E_HR_appendix}
\end{equation}
with boundary flux condition
\begin{equation}
\dot E(t)=F(0,t).
\label{eq:dEdt_flux_appendix}
\end{equation}
At the entrance boundary, the Marshak condition is
\begin{equation}
\sigma_{\mathrm{SB}}T_D^4(t)=\sigma_{\mathrm{SB}}T_s^4(t)+\frac{1}{2}F(0,t).
\label{eq:marshak_bc_appendix}
\end{equation}
Equations \eqref{eq:xF_HR_appendix}--\eqref{eq:marshak_bc_appendix} close the 1D problem.

\subsection{Implicit time-marching form for $T_s(t)$}
Following the Appendix-A construction used in the long manuscript, combine Eqs.~\eqref{eq:xF_HR_appendix} and \eqref{eq:E_HR_appendix} to get
\begin{equation}
E^2(t)=Z_1\,H^{\varepsilon}(t)\int_0^t H(t')\,dt',
\qquad
Z_1\equiv f^2\rho^{2(1-\mu)}(2+\varepsilon)(1-\varepsilon)C.
\label{eq:E2_identity_appendix}
\end{equation}
At time step $t_n=t_{n-1}+\Delta t$, define $H_n\equiv H(t_n)$ and
\begin{equation}
I_{n-1}\equiv\int_0^{t_{n-1}}H(t')\,dt'.
\end{equation}
Using a trapezoid approximation for the last interval,
\begin{equation}
\int_{t_{n-1}}^{t_n}H(t')\,dt'\approx\frac{H_{n-1}+H_n}{2}\Delta t,
\end{equation}
Eq.~\eqref{eq:E2_identity_appendix} becomes one nonlinear equation for $H_n$:
\begin{equation}
Z_1\left(I_{n-1}+\frac{H_{n-1}+H_n}{2}\Delta t\right)H_n^{\varepsilon}-E_n^2=0.
\label{eq:implicit_Hn_appendix}
\end{equation}
The input energy is advanced from the Marshak flux:
\begin{equation}
F_n=2\sigma_{\mathrm{SB}}\left(T_{D,n}^4-T_{s,n-1}^4\right),
\qquad
E_n=E_{n-1}+\frac{F_n+F_{n-1}}{2}\Delta t.
\label{eq:Fn_En_appendix}
\end{equation}
Then solve Eq.~\eqref{eq:implicit_Hn_appendix} for the positive root $H_n>0$, recover
\begin{equation}
T_{s,n}=H_n^{1/(4+\alpha)},
\end{equation}
and update
\begin{equation}
I_n=I_{n-1}+\frac{H_{n-1}+H_n}{2}\Delta t,
\qquad
x_{F,n}^2=\frac{2+\varepsilon}{1-\varepsilon}C\,H_n^{-\varepsilon}I_n.
\end{equation}
This is the practical algorithm used to compute the 1D centerline front with Marshak boundary forcing.

\section{Wall-Energy Loss: Full Integral Form}
Let $t_0(x)$ be the arrival time of the front at axial location $x$, defined by
\begin{equation}
x_F\bigl(t_0(x)\bigr)=x.
\end{equation}
After arrival, the local wall element is exposed for $\Delta t_x=t-t_0(x)$. For wall material $m$ (Au/Cu/Be), the absorbed areal energy is written in self-similar form
\begin{equation}
E_{w,m}(\Delta t_x;T_s)\approx A_m\,T_s^{p_m}(t_0)\,\Delta t_x^{q_m},
\qquad \Delta t_x>0.
\label{eq:wall_areal_generic_appendix}
\end{equation}
For the Au case used in the paper, one common fit is
\begin{equation}
E_{w,\mathrm{Au}}(\Delta t_x)\simeq0.59\,T_0^{3.35}\,\Delta t_x^{0.59}
\quad [\mathrm{hJ/mm^2}],
\label{eq:wall_Au_fit_appendix}
\end{equation}
with coefficients from the cited subsonic wall-heating model.

The total wall-loss energy is then
\begin{equation}
E_W(t)=2\pi\int_0^{x_F(t)}R(x,t)\,E_{w,m}\!\left(t-t_0(x)\right)dx.
\label{eq:EW_total_appendix}
\end{equation}
In differential form, this is equivalent to
\begin{equation}
\dot E_W(t)=2\pi\int_0^{x_F(t)}R(x,t)\,\dot E_{w,m}\!\left(t-t_0(x)\right)dx.
\label{eq:dEWdt_total_appendix}
\end{equation}
Numerically, Eqs.~\eqref{eq:EW_total_appendix}--\eqref{eq:dEWdt_total_appendix} are evaluated on the same front-time grid used for $x_F(t)$, with interpolation for $t_0(x)$.

\section{Wall Ablation and Density Feedback}
For subsonic wall response (relevant for Au and Cu sleeves in our regime), the inward wall velocity is modeled as
\begin{equation}
u_w(\Delta t_x;T_s)\approx B_m\,\tilde u(\rho_w)\,T_s^{\eta_m}(t_0)\,\Delta t_x^{\xi_m}.
\label{eq:uw_generic_appendix}
\end{equation}
For Au, a representative fit is
\begin{equation}
u_{w,\mathrm{Au}}\simeq510.1\,\tilde u(\rho_{\mathrm{Au}})\,T_0^{0.716}\,\Delta t_x^{0.036}
\quad [\mathrm{km/s}].
\end{equation}
The effective tube radius follows
\begin{equation}
R(x,t)=R_0-\int_{t_0(x)}^t u_w\!\left(t'-t_0(x)\right)dt'.
\label{eq:R_xt_appendix}
\end{equation}
The foam volume inside the heated region is
\begin{equation}
V(t)=\int_0^{x_F(t)}\pi R^2(x,t)\,dx,
\end{equation}
and the mean density update is
\begin{equation}
\rho(t)=\rho_0\frac{V_0}{V(t)}.
\label{eq:rho_update_appendix}
\end{equation}
Because $C\propto\rho^{2-\mu-\lambda}$, this introduces a density-history correction in the effective front coefficient (as used in the implementation).

\section{Radial Eigenvalue Problem and Albedo Coupling}


\section{Simulation details}
The simulation introduced in the letter is a 2D diffusion simulation without hydrodynamics, using cylindral geometry. The equations discritization is done through finite diffrencing on the following coupled equations:
\begin{equation}
\frac{\partial E}{\partial t}
=
\frac{1}{r}\frac{\partial}{\partial r}\left(rD\frac{\partial E}{\partial r}\right) + \frac{\partial}{\partial z}\left(D\frac{\partial E}{\partial z}\right)
-
\chi c\sigma(E - U).
\label{eq:E_continuous_conservative}
\end{equation}
This is the form we use throughout the derivation below.

\subsection{Material Energy Equation}

The material energy density $U(r,z,t)$ evolves according to:
\begin{equation}
\frac{1}{\beta}\,\frac{\partial U}{\partial t}
=
\chi c\sigma(E - U),
\label{eq:U_continuous}
\end{equation}
where $\beta = \frac{\partial U_R}{\partial U_m}$ is the ratio of radiation to material heat capacity (dimensionless coupling strength).

\subsection{Decomposition of the Diffusion Operator}

Define the space-dependent diffusion coefficient:
\begin{equation}
D(r,z,t)
\equiv
\frac{c}{3\sigma(T(r,z,t))}.
\end{equation}

\end{document}
